\documentclass[10pt]{beamer}

\usetheme{metropolis}
\definecolor{WiLabRed}{RGB}{197,18,48}
\setbeamercolor{frametitle}{fg=white,bg=WiLabRed}
\setbeamercolor{progress bar}{fg=WiLabRed!90}
\setbeamercolor{title separator}{fg=WiLabRed!90}
\setbeamercolor{progress bar in section page}{fg=WiLabRed!90}
\setbeamercolor{background canvas}{bg=white}
\setbeamercolor{alerted text}{fg=WiLabRed!90}

\usepackage{appendixnumberbeamer}

\usepackage{booktabs}
\usepackage[scale=2]{ccicons}

\usepackage{pgfplots}
\usepgfplotslibrary{dateplot}

\usepackage{xspace}
\newcommand{\themename}{\textbf{\textsc{metropolis}}\xspace}

\usepackage{marvosym}
%\usepackage{subfig}
\usepackage{graphicx}\graphicspath{{images/}}
\usepackage{subcaption}
\usepackage[framed]{./support/mcode}
\usepackage{listings}

\title{Lecture 10: Signal Space Concept}
\subtitle{\textit{Software Defined Radio for Engineers} (Collins~\textit{et al.}), \textsection{4.3}}
\date{}
\author{\textbf{Alexander M. Wyglinski, Ph.D.}}
\institute{ \vspace*{1in}\hfill\includegraphics[height=1.125cm]{wilab_logo-A70916.eps} \qquad \includegraphics[height=1.125cm]{WPI_Inst_Prim_FulClr.eps}}
% \titlegraphic{\hfill\includegraphics[height=1.5cm]{logo.pdf}}

% Foot for all slides
\setbeamertemplate{frame footer}{\tiny \copyright~2018 by Alexander Wyglinski. This work is licensed under the Creative Commons Attribution-ShareAlike 4.0 International License. To view a copy of this license, visit http://creativecommons.org/licenses/by-sa/4.0/.}

\begin{document}

%\captionsetup[subfigure]{labelformat=empty}

%%%%%%%%%%%%%%%%%%%%%%%%%%%%%%%%%%%%%%%%%%%%%%%%%%%%%%%%%%

\maketitle




%%%%%%%%%%%%%%%%%%%%%%%%%%%%%%%%%%%%%%%%%%%%%%%%%%%%%%%%%%

\frame
{
  \frametitle{Signal Vectors}

    \begin{itemize}
        \item Let $\phi_j(t)$ be an orthnormal set of functions on the time interval $[0,T]$ such that:
        \begin{equation}
            \int\limits_0^T\phi_i(t)\phi_j(t)dt=\left\{
            \begin{array}{ll}
              1 & i=j \\
              0 & \mathsf{otherwise}
            \end{array}
            \right.\nonumber
        \end{equation}
        \item Let $s_i(t)$ be the modulation signal that we can represent in terms of the orthonormal functions:
        \begin{equation}
            s_i(t)=\sum\limits_{k=1}^{N}s_{ik}\phi_k(t)
        \end{equation}
        which can be represented by the vector:
        \begin{equation}
            s_i(t)\rightarrow\mathbf{s}_i=(s_{i1},~s_{i2},~s_{i3},\ldots{s_{iN}})
        \end{equation}
    \end{itemize}

}
\frame
{
  \frametitle{Vector Representation of $s_i(t)$}
  
\begin{itemize}
 \item Let us look at a sample vector representation of $s_i(t)$ in three dimensional space using basis functions $\phi_1(t)$, $\phi_2(t)$, and $\phi_3(t)$
\end{itemize}

}
\frame
{
  \frametitle{Vector Manipulations}

    \begin{itemize}
        \item To find $s_{ik}$, solve:
        \begin{equation}
            \int\limits_0^Ts_i(t)\phi_l(t)dt=\sum\limits_{k=1}^{N}s_{ik}\int\limits_0^T\phi_k(t)\phi_l(t)dt=s_{il}
        \end{equation}
        \item The vector dot product between $s_i(t)$ and $s_j(t)$ is equal to:
        \begin{equation}
            \int\limits_0^Ts_i(t)s_j(t)dt=\mathbf{s}_i\cdot\mathbf{s}_j=\rho_{ij}
        \end{equation}
        while the energy of a signal $s_i(t)$ is equal to:
        \begin{equation}
            E_{s_i}=\int\limits_0^Ts_i^2(t)dt=\mathbf{s}_i\cdot\mathbf{s}_i=||\mathbf{s}_i||^2
        \end{equation}
    \end{itemize}
}
\frame
{
  \frametitle{Euclidean Distance}

    \begin{itemize}
        \item To compute the {\it Euclidean Distance} using signal space vectors, we need to solve:
        \begin{equation}
        \begin{split}
            d_{\min}^2&=\int\limits_0^T\Delta{s_{ij}^2(t)}dt=\int\limits_0^T(s_i(t)-s_j(t))^2dt\\
            &=||\mathbf{s}_i-\mathbf{s}_j||^2=(\mathbf{s}_i-\mathbf{s}_j)\cdot(\mathbf{s}_i-\mathbf{s}_j)\\
            &=E_{s_i}+E_{s_j}-2\rho_{ij}\nonumber
        \end{split}
        \end{equation}
        where:
        \begin{equation}
            \rho_{ij}=\int\limits_0^Ts_i(t)s_j(t)dt=\mathbf{s}_i\cdot\mathbf{s}_j
        \end{equation}
    \end{itemize}
}
\frame
{
  \frametitle{Solving for the Power Efficiency}

    \begin{itemize}
        \item Choose orthonormal basis functions $\phi_i(t)$, $i=1,2,\ldots,k$, where $k$ is the dimension of the signal space
        \item Find $\mathbf{s}_i$, $i=1,2,\ldots,M$ where $\mathbf{s}_i=(s_{i1},~s_{i2},\ldots{s_{ik}})$ and $s_{ij}=\int\limits_0^Ts_i(t)\phi_j(t)dt$
        \item Consequently, solve for:
        \begin{equation}
        \begin{split}
            d_{\min}^2&=\min\limits_{i\neq{j}}||{\mathbf{s}_i-\mathbf{s}_j}||^2\\
            \bar{E}_s&=\frac{1}{M}\sum\limits_{i=1}^{M}||\mathbf{s}_i||^2\\
            \bar{E}_b&=\bar{E}_s/\log_2(M)\\
            \varepsilon_p&=d_{\min}^2/\bar{E}_b\nonumber
        \end{split}
        \end{equation}
    \end{itemize}
}



\end{document}
